\section{Анализ предметной области}
\subsection{Управление проектами как область деятельности}

ГОСТ Р ИСО 21500--2023 определяет проект как временное предприятие, направленное на создание уникального продукта, услуги или результата \cite{gost21500}. ГОСТ Р 54869--2011 трактует проект иначе -- как комплекс взаимосвязанных мероприятий, направленный на создание уникального результата в условиях ограничений по времени и ресурсам \cite{gost54869}.

Из этих определений видны четыре признака, по которым проект отличается от операционной деятельности:
\begin{enumerate}
    \item Конечные сроки -- у проекта есть дата начала и дата, к которой результат должен быть готов.
    \item Уникальность результата -- речь не о том, что подобного никогда не делалось. Имеется в виду уникальное сочетание контекста, требований и условий: два внешне одинаковых по ТЗ сайта, сделанных для разных компаний, все равно остаются разными проектами \cite{mazur2014}.
    \item Ресурсные ограничения -- любой проект тратит время людей, бюджет и оборудование. Балансировка между этими ограничениями и объемом работ -- большая часть работы руководителя.
    \item Измеримая цель -- без сформулированной цели не понять, в какой момент проект можно считать завершенным.
\end{enumerate}

Управление проектом, по PMBOK, -- это применение знаний, навыков, инструментов и методов к работам проекта для удовлетворения требований, которые к нему предъявлены \cite{pmbok6}. Руководитель в любой момент должен знать состав работ, исполнителей и сроки. Без централизованного хранения эти сведения быстро расходятся по почте и переписке, теряют актуальность.

PMBOK делит процессы управления на пять групп: инициация, планирование, исполнение, мониторинг и контроль, завершение. Первые четыре образуют основной поток работ от запуска до сдачи проекта. Мониторинг и контроль идут параллельно планированию и исполнению -- они отвечают за то, чтобы план вовремя корректировался по фактическим результатам. Взаимосвязь процессов показана на рисунке \ref{fig:pmbok}.

\begin{figure}[H]
\centering
\includegraphics[width=\textwidth]{pmbok}
\caption{Группы процессов управления проектом по PMBOK}
\label{fig:pmbok}
\end{figure}

ГОСТ Р 54869--2011 выделяет в проектной деятельности четыре ключевые роли: заказчик проекта, куратор, руководитель и команда проекта \cite{gost54869}. Основная часть рабочих процессов строится вокруг руководителя и команды -- именно их повседневные потребности определяют, что должна уметь система автоматизации.

Как самостоятельная дисциплина управление проектами оформилось во второй половине XX века. Сначала появились формальные методы планирования: диаграмма Ганта в 1910 году, методы PERT и CPM в конце 1950-х. В 1969 году учрежден Project Management Institute (PMI), в 1996 году вышло первое издание PMBOK -- Свода знаний по управлению проектами. К 1990-м годам управление проектами вышло за пределы инженерных и оборонных отраслей \cite{pmbok6, mazur2014}. Сегодня его принципы применяют в ИТ, маркетинге, образовании, государственных программах.

Чем больше проектов и распределенных команд, тем тяжелее держать планирование на бумаге и в почте. Отсюда возникла потребность в специализированных программных средствах.

\subsection{Методологии управления проектами}

Методология управления проектами -- это набор принципов и практик, который описывает, как проект планируется, исполняется и контролируется \cite{terehova_review}. Конкретные инструменты она не предписывает. Внутри методологии команда сама выбирает программы и форматы документов под свои задачи.

Условно такие методологии делят на две большие группы: классические (последовательные) и гибкие. Их принципиальное отличие показано на рисунке \ref{fig:waterfall_iter}.

\begin{figure}[H]
\centering
\includegraphics[width=\textwidth]{waterfall_iter}
\caption{Сравнение классической и гибкой моделей разработки}
\label{fig:waterfall_iter}
\end{figure}

В классической модели все этапы выстраиваются в одну цепочку: каждый следующий начинается только после завершения предыдущего. Гибкий подход устроен иначе -- работа разбивается на короткие повторяющиеся циклы, каждый из которых проходит все стадии и завершается работающим инкрементом продукта. На практике большинство команд комбинирует элементы обоих подходов -- такие комбинации называют гибридными.

\subsubsection{Классический подход}

Классическая модель опирается на представление о проекте как о цепочке последовательных стадий: формирование требований, проектирование, реализация, тестирование, ввод в эксплуатацию. Каждая следующая стадия начинается после завершения предыдущей. Такой подход известен как каскадная модель и впервые был сформулирован в 1970-х годах применительно к разработке программного обеспечения \cite{mazur2014}.

Каскадные методы применимы там, где требования зафиксированы на старте и не меняются по ходу работ: строительство, государственные контракты, разработка отраслевых регламентов. При соблюдении этого условия проект остается предсказуемым по срокам и стоимости.

Классический подход дал начало формальным инструментам сетевого планирования: диаграмме Ганта, методу критического пути и методу PERT \cite{pmbok6}. Все три до сих пор применяются в системах управления проектами, когда нужно учесть порядок выполнения работ и влияние задержки отдельной задачи на общий срок проекта.

\subsubsection{Диаграмма Ганта}

Диаграмма Ганта разработана американским инженером Генри Гантом в начале XX века и с тех пор остается стандартным инструментом планирования работ \cite{art_gantt}. Каждая задача показана горизонтальным отрезком, длина которого передает ее продолжительность. Левый край отрезка соответствует дате начала задачи, правый -- дате завершения. Когда между задачами добавляют связи, диаграмма позволяет определить критический путь -- самую длинную взаимосвязанную цепочку работ, которая задает нижнюю границу общей длительности проекта. Главное достоинство такого представления -- наглядность. На одном экране умещается план из десятков работ, и сразу видно, какие из них идут параллельно, а какие последовательно.

Диаграмма Ганта не является атрибутом исключительно классического подхода. В большинстве современных ИС~УП она реализована как базовое представление данных, доступное в том числе командам, работающим по гибким методологиям: задачи группируются по спринтам, зависимости отражают связи между инкрементами. Ее присутствие в каждой из рассматриваемых далее систем подтверждает, что это рабочий инструмент независимо от выбранной методологии.

\subsubsection{Agile и его частные методы}

Гибкие методологии (Agile) появились в ИТ как ответ на жесткость каскадной модели. В 2001 году группа разработчиков сформулировала <<Манифест гибкой разработки программного обеспечения>>: работающий продукт ставится выше исчерпывающей документации, сотрудничество с заказчиком -- выше контрактных формальностей, готовность к изменениям -- выше следования первоначальному плану \cite{art_agile_modern}.

Среди конкретных методов Agile наибольшее применение в ИТ получили Scrum и Kanban.

Scrum строит работу команды короткими циклами фиксированной длительности -- спринтами, обычно по две-четыре недели. На старте спринта команда отбирает задачи из общего журнала, в конце -- показывает работающий результат и проводит ретроспективу \cite{sutherland_scrum}. Ролей в команде три: владелец продукта, скрам-мастер и разработчики. Подход хорошо ложится на продуктовую разработку с коротким горизонтом планирования.

Kanban пришел в ИТ из производственной системы Toyota. Дэвид Андерсон в середине 2000-х адаптировал его под работу со знаниями \cite{anderson_kanban}. В основе -- доска с колонками по этапам и ограничение на число задач в работе одновременно. Жестких итераций тут нет, поэтому Kanban удобен там, где задачи идут потоком: поддержка, сопровождение, разбор инцидентов.

Экстремальное программирование (XP) больше про инженерные практики, чем про управление: парное программирование, разработка через тестирование, непрерывная интеграция, простота проектных решений \cite{beck_xp}. По части проектного управления XP менее формален, чем Scrum, но его инженерные приемы перешли в другие методологии.

\subsubsection{Гибридные подходы}

Большинство команд на практике не следует одной методологии строго -- они заимствуют элементы из разных подходов: итеративную работу из Scrum, доску с колонками из Kanban, формальное закрытие фаз из каскадной модели \cite{art_hybrid_pm}. Такие комбинации принято называть гибридными.

Самый известный гибридный формат -- Scrumban. От Scrum он берет ритм ретроспектив, а жесткие границы спринта меняет на непрерывный поток: задачи поступают на доску по мере необходимости, а лимит на число одновременных задач заимствуется из Kanban. Такая схема выручает команды со смешанной нагрузкой, где плановые задачи на новую функциональность чередуются с дефектами и запросами на поддержку.

Для крупных компаний с несколькими командами разработаны фреймворки масштабирования. SAFe (Scaled Agile Framework, масштабированный гибкий фреймворк) надстраивает над спринтами несколько уровней координации, где команды синхронизируются через квартальные программные инкременты, а на уровне портфеля остается стратегическое планирование, характерное для классической модели. За счет такой надстройки несколько команд могут работать в едином ритме, не теряя гибкости отдельного спринта.

Шестое издание PMBOK прямо признает: реальные проекты обычно идут по адаптированному процессу, а не по строгой реализации одной методологии \cite{pmbok6}. Для ИС~УП это означает простое требование. Система должна одновременно поддерживать спринтовое планирование, непрерывный поток задач и классические зависимости между работами, не заставляя команду выбирать что-то одно.

Ключевые характеристики рассмотренных методологий сведены в таблицу~\ref{compare:methods}.

\renewcommand{\arraystretch}{1.0}
\begin{xltabular}{\textwidth}{|>{\hspace{0pt}\raggedright\arraybackslash}p{3.3cm}|>{\hspace{0pt}\centering\arraybackslash}X|>{\hspace{0pt}\centering\arraybackslash}X|>{\hspace{0pt}\centering\arraybackslash}X|>{\hspace{0pt}\centering\arraybackslash}X|}
\caption{Сравнение основных методологий управления проектами\label{compare:methods}}\\ \hline
\centrow Признак & \centrow Каскадная & \centrow Scrum & \centrow Kanban & \centrow XP \\ \hline
\endfirsthead
\continuecaption{Продолжение таблицы \ref{compare:methods}}
\centrow Признак & \centrow Каскадная & \centrow Scrum & \centrow Kanban & \centrow XP \\ \hline
\finishhead
Тип выполнения & Последовательный & Итеративный & Непрерывный & Итеративный \\ \hline
Цикл планирования & Весь проект единовременно & Спринт 1--4~нед. & Нет фиксир. цикла & Итерация 1--2~нед. \\ \hline
Реакция на изменения & Затруднена & Возможна между спринтами & В любой момент & Приветствуется \\ \hline
Документирование & Подробное, заранее & Минимальное & Минимальное & Минимальное \\ \hline
Область применения & Строительство, госзаказ & Продуктовая разработка & Поддержка, сопровождение & Разработка ПО \\
\end{xltabular}

\subsection{Информационные системы для управления проектами}

Информационная система управления проектами -- это программный комплекс, который автоматизирует планирование, исполнение и контроль работ по одному или нескольким проектам \cite{safonova_is}. Данные о проектах, задачах, ресурсах и сроках хранятся в одной базе, и команда работает с единой версией плана.

Первые специализированные программы для управления проектами появились в 1970-х. Работали они на мейнфреймах и обслуживали крупные промышленные и строительные корпорации. В 1984 году вышел Microsoft Project для MS-DOS -- один из первых массовых настольных продуктов в этой области \cite{akhmadullin_pms}. Когда персональные компьютеры стали повсеместными, рынок расширился за пределы крупных организаций. В 2000-е планирование переехало в браузер -- чтобы начать работу, достаточно было зарегистрироваться в сервисе. Сегодня рынок делят между собой корпоративные платформы, специализированные таск-трекеры и отраслевые решения \cite{akhmadullin_pms}.

\subsubsection{Ключевые сущности предметной области}

Продукты на рынке отличаются интерфейсом и набором надстроек, но базовый состав сущностей у большинства ИС~УП один и тот же. Главная -- проект. Он задает контекст работ, сроки, цель и состав участников. Без проекта система превращается в плоский список задач, где трудно понять, к какой инициативе относится та или иная работа.

Проект разбивается на задачи. У каждой задачи есть описание, даты начала и завершения, продолжительность и место в общей последовательности. Задача в проектной системе не существует сама по себе -- она часть календарного плана.

Состав команды неоднороден. Один участник заводит проект и собирает остальных, у другого основная работа -- правка задач, третий просто следит за ходом дел и не должен иметь возможности что-то изменить. Если система не разграничивает эти случаи, пользователи получают либо лишние права, либо избыточные ограничения. Поэтому ИС~УП хранит не только сам факт участия пользователя в проекте, но и его роль.

Зависимости между задачами фиксируют порядок выполнения работ. Без них в плане легко получить противоречие: задача-преемник начинается раньше предшественника. В небольшом проекте из пяти-шести задач этим можно пренебречь, но для больших планов связи -- единственный способ удержать управляемость.

\subsubsection{Классификация систем управления проектами}

По масштабу ИС~УП обычно делят на три группы \cite{safonova_is}. На одном краю спектра -- корпоративные системы, в которых сотни пользователей и есть интеграция с CRM, ERP, учетом рабочего времени. На другом краю -- индивидуальные планировщики, которыми пользуется один человек. Промежуточный сегмент занимают системы для малых и средних команд, без тяжелых интеграций, но и без ограничения «один пользователь».

Вторая ось -- модель распространения. Облачные системы (SaaS, программное обеспечение как услуга) размещаются на инфраструктуре поставщика, обновления и резервные копии тоже на его стороне. Пользователю достаточно браузера. Локальные решения устанавливаются на сервер заказчика, и эксплуатация переходит к нему. Большую долю рынка занимают облачные продукты: подключиться проще, входная стоимость ниже. Локальная установка нужна там, где требования к конфиденциальности не позволяют отдавать данные стороннему поставщику \cite{akhmadullin_pms}.

Третий признак -- специализация. Универсальные системы подходят для любого проекта и дают гибкие настройки процессов. Отраслевые поставляются с готовыми шаблонами под конкретную сферу: разработку ПО, строительство, маркетинг, научные исследования.

\subsubsection{Типовая функциональность}

Подходы и масштабы у систем разные, но базовый набор функций совпадает у большинства ИС~УП:
\begin{itemize}
    \item реестр проектов с описанием, сроками и ответственными;
    \item декомпозиция проекта на задачи -- каждой можно назначить исполнителя, срок, статус и приоритет;
    \item диаграмма Ганта или иная визуализация календарного плана;
    \item зависимости между задачами для отражения порядка работ;
    \item разграничение прав доступа между участниками проекта;
    \item история изменений и операций;
    \item комментарии в контексте задачи;
    \item отчеты о состоянии проекта.
\end{itemize}

Поверх этого ядра конкретные продукты добавляют то, что считают нужным для своей аудитории. В одних системах это учет рабочего времени и автоматизация повторяющихся действий, в других -- связь с системами контроля версий и мобильные клиенты. Аналитические панели с метриками встречаются ближе к корпоративному сегменту.

\subsubsection{Проблемы применения ИС~УП в небольших командах}

Крупные корпоративные платформы покрывают широкий круг задач, но требуют времени на настройку, обучение пользователей и администрирование. Для команды из пяти человек это явный перебор: вместо ведения проекта она обслуживает инструмент.

С простыми досками задач история обратная. На старте все удобно, но как только в проекте появляются сроки, взаимосвязанные работы и разные уровни доступа, досок перестает хватать. Перемещать карточки между колонками можно, а целостной картины сроков нет. Руководитель отдельно ведет календарный план в Excel, вручную проверяет зависимости и объясняет участникам, почему одну задачу нельзя начать раньше другой.

Поэтому небольшой команде нужна система, которая запускается без длительного внедрения и при этом покрывает базовый проектный контроль -- единый перечень проектов, календарь задач, зависимости и разграничение прав. Без любого из этих пунктов инструмент либо перегружает команду настройкой, либо перестает справляться по мере роста проекта.

\subsubsection{Технологические тенденции}

За последние годы рынок ИС~УП заметно сдвинулся по нескольким направлениям.

Первое и самое заметное -- переход к веб-архитектуре с тонким клиентом. Большинство систем стало браузерными, и пользователь работает с любого устройства, а обновления поставщик выпускает централизованно. Норма распространилась даже на исторически настольные продукты: у Microsoft Project появилась веб-версия Project for the web.

Открытый программный интерфейс из конкурентного преимущества превратился в обязательную часть. REST API сегодня есть почти у всех систем. Через него систему встраивают во внутреннюю инфраструктуру предприятия, выгружают данные в аналитические панели, связывают с системами контроля версий. Продукты, ориентированные на программные команды, идут дальше и добавляют готовые связки с Git, GitHub, GitLab и сервисами непрерывной интеграции -- так привязка коммитов к задачам выполняется автоматически, а ход работ отслеживается по реальным артефактам, а не только по статусам в отчетах.

Мобильные клиенты подтягиваются к веб-версии и по функциональности уже почти не уступают ей. Поддерживается работа без сети с последующей синхронизацией. Всплывающие уведомления особенно полезны в распределенных командах и проектах с выездными работами.

Отдельная тенденция -- встроенная аналитика. В системы стали добавлять панели с проектными метриками (диаграммы выгорания, среднее время выполнения задачи, пропускная способность команды, прогноз сроков), часть которых опирается на машинное обучение: система ищет зависшие задачи и оценивает риски срыва сроков.

\subsection{Сравнительный анализ существующих решений}

Для сопоставления выбраны восемь систем -- три зарубежные с широкой базой пользователей в ИТ-среде и пять решений, доступных в Российской Федерации без ограничений и активно развивающихся в последние годы \cite{yougile_review}.

\subsubsection{Краткая характеристика рассматриваемых систем}

\textbf{Jira} -- одна из самых распространенных систем для управления ИТ-проектами. Поддерживает Scrum- и Kanban-доски, позволяет настраивать рабочие процессы, у нее большая экосистема плагинов. Выпускается компанией Atlassian и считается отраслевым стандартом отслеживания задач в командах разработки ПО. В марте 2022 года Atlassian приостановил продажи в РФ и Беларуси, к октябрю 2022 года компания полностью свернула здесь свою деятельность. В августе 2023 года Atlassian уведомил пользователей с российскими и белорусскими аккаунтами о принудительном отключении облачных сервисов в течение 30 дней. С 15 февраля 2024 года серверные версии Atlassian остались без официальной поддержки, обновлений безопасности и исправлений ошибок. Для отечественных команд Jira после этого фактически непригодна.

\textbf{Trello} -- простая Kanban-доска от того же разработчика. Низкий порог входа, большое сообщество. В 2017 году продукт перешел в собственность Atlassian. Платная функциональность скромная, поэтому Trello удобна для одного человека или небольшой группы. Для проектов, где нужны зависимости задач и диаграмма Ганта, инструмент не подходит.

\textbf{Asana} -- универсальная система для команд среднего и крупного размера. Поддерживает разные представления одних и тех же данных: список, доска, диаграмма Ганта, календарь. Ориентирована на кросс-функциональные команды -- маркетинг, операционные группы, проектный офис, разработка ПО проработана слабее. Русского интерфейса нет. С ноября 2025 года бесплатный тариф Asana Personal ограничили двумя пользователями, ранее зарегистрированные команды сохраняют прежний лимит в десять человек.

\textbf{YouTrack} -- разработка JetBrains, рассчитанная на команды разработки. Особенности -- запросы на естественном языке для поиска задач и тесная связка с другими продуктами JetBrains. При локальном развертывании бесплатна для команд до десяти пользователей. Это удобно небольшим техническим коллективам, которые хотят держать данные у себя. Облачная версия тоже есть.

\textbf{Битрикс24} -- комплексная российская платформа, объединяющая управление проектами с CRM, документооборотом и корпоративными коммуникациями \cite{yougile_review}. Продукт компании 1С-Битрикс поставляется в облачном и локальном вариантах. Бесплатный тариф не ограничивает число пользователей, но дает только 5~ГБ облачного хранилища и ограничивает функциональность CRM и автоматизаций. Подходит крупным компаниям, но порог входа высокий -- модулей много, и управление задачами теряется среди них.

\textbf{Kaiten} -- российская система с упором на Kanban и Scrum. Сильна визуализацией потока работ и аналитикой, целевая аудитория -- команды разработки и продуктовые отделы. Из особенностей -- встроенная аналитика потока: среднее время выполнения задач, диаграммы накопленного потока, пропускная способность команды за произвольный период.

\textbf{YouGile} -- российская система. Объединяет Kanban-доску со встроенным мессенджером прямо в карточке задачи. Переписку по конкретной работе ведут не в отдельных чатах, а в карточке -- к ней потом всегда можно вернуться. Бесплатна для команд до десяти человек, порог входа низкий.

\textbf{Microsoft Project} -- классический инструмент с глубокой поддержкой сетевого планирования и расчета критического пути. Первая версия вышла в 1984 году. Сегодня продукт идет в настольном варианте Project Professional и облачном Project for the web. По сетевому планированию, управлению ресурсами и межпроектным связям он богаче большинства конкурентов, но и порог освоения, и цена лицензии заметно выше. Применяется в проектах с жестким каскадным управлением, в основном вне ИТ.

\subsubsection{Сравнение по ключевым критериям}

Для сопоставления выбраны критерии, существенные для команды, которая ведет несколько параллельных программных проектов с распределенными участниками. Один из критериев требует пояснения. Под четырьмя типами зависимостей задач понимаются классические связи сетевого планирования: <<финиш-старт>>, <<старт-старт>>, <<финиш-финиш>> и <<старт-финиш>>. Их графическое представление приведено на рисунке \ref{fig:dep_types}.

\begin{figure}[H]
\centering
\includegraphics[width=\textwidth]{dep_types}
\caption{Четыре типа зависимостей задач в проекте}
\label{fig:dep_types}
\end{figure}

В таблице~\ref{compare:availability} собраны характеристики доступности: модель распространения, наличие русского интерфейса, бесплатного тарифа и фактический доступ из Российской Федерации. В таблице~\ref{compare:features} -- ключевые функциональные возможности: диаграмма Ганта, связи задач, открытый программный интерфейс, мобильное приложение.

\renewcommand{\arraystretch}{1.0}
\begin{xltabular}{\textwidth}{|>{\raggedright\arraybackslash}p{2.4cm}|T|T|T|T|}
\caption{Сравнение систем по доступности и условиям использования\label{compare:availability}}\\ \hline
\centrow Система & \centrow Модель распро\-странения & \centrow Русский язык & \centrow Бесплат\-ный тариф & \centrow Доступ из РФ \\ \hline
\endfirsthead
\continuecaption{Продолжение таблицы \ref{compare:availability}}
\centrow Система & \centrow Модель распро\-странения & \centrow Русский язык & \centrow Бесплат\-ный тариф & \centrow Доступ из РФ \\ \hline
\finishhead
Jira & Обе & Частично & 10 чел. & Огранич. \\ \hline
Trello & Облачная & Да & 10 досок & Да \\ \hline
Asana & Облачная & Нет & 2 чел. & Огранич. \\ \hline
YouTrack & Обе & Да & 10 чел. & Да \\ \hline
Битрикс24 & Обе & Да & Без огр. & Да \\ \hline
Kaiten & Обе & Да & 5 чел. & Да \\ \hline
YouGile & Обе & Да & 10 чел. & Да \\ \hline
MS Project & Обе & Да & Нет & Да
\end{xltabular}

\begin{xltabular}{\textwidth}{|>{\raggedright\arraybackslash}p{2.4cm}|T|T|T|T|}
\caption{Сравнение систем по функциональным возможностям\label{compare:features}}\\ \hline
\centrow Система & \centrow Диаграмма Ганта & \centrow Связи задач & \centrow Открытый API & \centrow Мобильное приложение \\ \hline
\endfirsthead
\continuecaption{Продолжение таблицы \ref{compare:features}}
\centrow Система & \centrow Диаграмма Ганта & \centrow Связи задач & \centrow Открытый API & \centrow Мобильное приложение \\ \hline
\finishhead
Jira & Частично (Timeline) & 4 типа & Да & Да \\ \hline
Trello & Через плагин & Нет & Да & Да \\ \hline
Asana & Да & 1 тип & Да & Да \\ \hline
YouTrack & Да & Да & Да & Да \\ \hline
Битрикс24 & Да & 4 типа & Да & Да \\ \hline
Kaiten & Да & Да & Да & Да \\ \hline
YouGile & Да & 1 тип & Да & Да \\ \hline
MS Project & Да & 4 типа & Да & Огранич.
\end{xltabular}

\subsubsection{Выводы по сравнению}

Зарубежные Jira и Asana из РФ доступны с ограничениями -- для постоянной работы российских команд они отпадают. Trello, MS Project и все рассмотренные российские системы доступны без оговорок.

Из доступных Asana русского интерфейса не имеет вовсе, у Jira локализация частичная. Команде, у которой английский не основной рабочий язык, это лишний барьер.

Trello -- самая простая и популярная из зарубежных, но зависимости между задачами не поддерживает, нативной диаграммы Ганта тоже нет (только через сторонние плагины). Для проектов, где план строится на связях работ и сетевых графиках, не подходит.

Все четыре типа связей сетевого планирования среди свободно доступных в РФ систем поддерживают только Битрикс24 и MS Project. Kaiten и YouTrack дают зависимости без разделения на типы, YouGile -- только <<финиш-старт>>.

У MS Project нет бесплатного тарифа -- для команд с минимальным бюджетом не вариант. Kaiten на бесплатном тарифе допускает пять пользователей, для среднего коллектива мало. Битрикс24 пользователей не ограничивает, но это универсальная корпоративная платформа с десятками модулей. Управление проектами там одна из многих подсистем, и порог входа высокий.

Получается, что ни одно из рассмотренных специализированных решений не объединяет в одной точке нативную диаграмму Ганта, четыре типа зависимостей, русский интерфейс, разумный бесплатный тариф и устойчивый доступ из РФ. Системы, закрывающие функциональную часть, либо уходят из России, либо превращаются в универсальную корпоративную платформу.

\subsection{Выводы по анализу предметной области}

Рабочий план -- это не плоский список задач. Чтобы он оставался актуальным, в системе должны храниться связи между задачами, сроки и состав участников. Когда эти сведения распределены по разным файлам и переписке, между ними появляется рассинхрон. План в Excel расходится со списком в трекере, ответственного за конкретную задачу никто не помнит, и руководитель тратит на сверку больше времени, чем на сами работы.

Веб-приложение, доступное всем участникам без установки клиента, удобно для распределенной работы. Руководитель и команда могут находиться в разных городах и при этом видеть один и тот же актуальный план.

Диаграмма Ганта удобна тем, что собирает в одно представление список задач, сроки и зависимости. Но она должна быть рабочим инструментом, а не картинкой в отчете. Любое изменение -- сдвиг дат, перестановка задач, новая связь -- обязано сразу попадать в данные проекта.

Прав доступа достаточно трех уровней. Полный набор прав -- только у создателя проекта: он распоряжается и параметрами, и составом команды. Промежуточный уровень нужен тем, кто пишет план и расставляет связи между задачами, но не должен иметь возможности удалить проект целиком или поменять состав участников. Наконец, тем, кто работает в проекте только для справки, достаточно доступа на чтение. Эти три уровня покрывают потребности небольшой команды и не требуют отдельного администрирования.

Разрабатываемая ИС~УП нацелена на конкретную практическую задачу -- календарное планирование с возможностью совместной работы. В рамках этой задачи решается ведение проектов и задач, отображение плана на диаграмме Ганта с зависимостями между работами и совместный доступ нескольких пользователей с разграничением их прав.
